% ============================================================
%  LOG DES ANOMALIES — Phase 3 : Exécution (recette)
%  Time'Eats — Cellule PMO
% ============================================================
\documentclass[11pt,a4paper,oneside]{report}
\usepackage{../style/consulting}
\usepackage{pdflscape}
\hypersetup{pdftitle={Log des Anomalies — Time'Eats PMO}}
\pagestyle{fancy}\fancyhf{}
\renewcommand{\headrulewidth}{0pt}
\fancyhead[L]{\begin{tikzpicture}[remember picture,overlay]
  \draw[cnNavy,line width=0.5pt](0,0)--(\textwidth,0);
\end{tikzpicture}\vspace{2pt}\timeeatslogo\hspace{4pt}\small\textcolor{cnSlate}{Time'Eats \textbf{·} Log des Anomalies — \textit{[Nom projet]}}}
\fancyhead[R]{\small\textcolor{cnSlate}{\thepage}}
\fancyfoot[C]{\tiny\textcolor{cnSlate}{Document confidentiel — Cellule PMO — CESI MAALSI 2026}}
\begin{document}\small

\begin{center}
{\LARGE\bfseries\color{cnNavy} Log des Anomalies}\\[4pt]
\textcolor{cnOrange}{\rule{10cm}{1pt}}\\[4pt]
{\large\color{cnSlate} Phase 3/5 — Exécution \& Recette \quad|\quad Projet : \textit{[Nom du projet]}}
\end{center}

\vspace{4pt}
\renewcommand{\arraystretch}{1.4}
\begin{tabularx}{\textwidth}{L{3.5cm} Y L{3.5cm} Y}
\toprule
\rowcolor{cnNavy}\th{Projet} & \th{\textit{[Nom]}} & \th{Dernière MAJ} & \th{\textit{[JJ/MM/AAAA]}} \\
\midrule
Responsable & \textit{[Prénom NOM]} & Version & \textit{[X.X]} \\
\bottomrule
\end{tabularx}

\vspace{4pt}
\begin{finding}
\textbf{Sévérités :}
\textbf{P1 Bloquant} — empêche la recette / la mise en prod. \quad
\textbf{P2 Majeur} — fonctionnalité altérée, contournement difficile. \quad
\textbf{P3 Mineur} — gêne utilisateur, contournement simple. \quad
\textbf{P4 Cosmétique} — affichage, libellé.
\end{finding}

\begin{landscape}
\renewcommand{\arraystretch}{1.4}
\begin{tabularx}{\linewidth}{C{1.2cm} C{2cm} L{3cm} C{1.5cm} L{3cm} L{3cm} C{2cm} C{2cm}}
\toprule
\rowcolor{cnNavy}
\th{ID} & \th{Date détection} & \th{Description} & \th{Sévérité} & \th{Module/Écran} & \th{Correction apportée} & \th{Responsable} & \th{Statut} \\
\midrule
ANO-001 & \textit{[JJ/MM]} & \textit{[Description courte]} & \textit{[P1]} & \textit{[Module]} & \textit{[Correction]} & \textit{[ESN]} & \textit{[Ouvert]} \\
\rowcolor{cnIce}
ANO-002 & \textit{[JJ/MM]} & \textit{[Description courte]} & \textit{[P2]} & \textit{[Module]} & \textit{[Correction]} & \textit{[ESN]} & \textit{[Résolu]} \\
ANO-003 & \textit{[JJ/MM]} & \textit{[Description courte]} & \textit{[P3]} & \textit{[Module]} & \textit{[Correction]} & \textit{[DSI]} & \textit{[Résolu]} \\
\rowcolor{cnIce}
ANO-004 & \textit{[JJ/MM]} & \textit{[Description courte]} & \textit{[P2]} & \textit{[Module]} & \textit{[En cours]} & \textit{[ESN]} & \textit{[En cours]} \\
ANO-005 & \textit{[JJ/MM]} & \textit{[Description courte]} & \textit{[P4]} & \textit{[Module]} & \textit{[Correction]} & \textit{[ESN]} & \textit{[Résolu]} \\
\bottomrule
\end{tabularx}
\end{landscape}

\section*{Tableau de synthèse}

\renewcommand{\arraystretch}{1.3}
\begin{tabularx}{\textwidth}{L{3cm} C{2cm} C{2cm} C{2cm} C{2cm}}
\toprule
\rowcolor{cnNavy}\th{Sévérité} & \th{Total détectés} & \th{Résolus} & \th{En cours} & \th{Ouverts} \\
\midrule
P1 Bloquant   & \textit{[X]} & \textit{[X]} & \textit{[X]} & \textit{[X]} \\
\rowcolor{cnIce}
P2 Majeur     & \textit{[X]} & \textit{[X]} & \textit{[X]} & \textit{[X]} \\
P3 Mineur     & \textit{[X]} & \textit{[X]} & \textit{[X]} & \textit{[X]} \\
\rowcolor{cnIce}
P4 Cosmétique & \textit{[X]} & \textit{[X]} & \textit{[X]} & \textit{[X]} \\
\midrule
\rowcolor{cnNavy!10}\textbf{TOTAL} & \textbf{\textit{[X]}} & \textbf{\textit{[X]}} & \textbf{\textit{[X]}} & \textbf{\textit{[X]}} \\
\bottomrule
\end{tabularx}

\begin{alert}
Seuil de GO / NO-GO pour la mise en production : \textbf{0 anomalie P1 ouverte} et \textbf{P2 ouverts $\leq$ 2} avec plan de correction \textless{} 5 jours ouvrés.
\end{alert}

\end{document}
